\chapter{Public Services and Human Professions}
\label{ch:15}

The preceding chapters of this part have considered education and care as domains in which the allocation of tasks between machines and persons is dictated by the nature of the activity. This chapter turns to the broader structure of public services, and to the professions through which they are delivered, and it takes up the question that most strongly affects the political reception of any program of automation, which is what happens to the persons who do the work. The question is not secondary. A program that proposes to increase the capacity of society to direct its own future cannot rest on the premise that a substantial part of society will lose its livelihood and its standing, and the history of technological transitions provides abundant reason to doubt the comfortable assurance that displaced workers will find equally good work elsewhere. The assurance has sometimes been borne out in the aggregate and has often failed for particular groups and regions, with consequences that lasted a generation.

An old observation in economics provides a useful starting point. Baumol and Bowen~\cite{baumol1966}, writing in 1966 about the performing arts, noted that a string quartet requires the same number of performers and the same amount of time to play a given piece today as in the nineteenth century, so that the productivity of the activity has not risen, while the wages of musicians have risen with wages in the rest of the economy, since otherwise musicians would leave for other occupations. The cost of live performance therefore rises relative to the cost of goods produced in sectors where productivity grows. The pattern, since called the cost disease, applies to any activity in which the labor is itself the product, which includes much of education, health care, and personal service. Its implication is that as the economy becomes more productive in manufacturing and information processing, the relative cost of services of this kind increases, and a growing share of income and of employment is devoted to them. The pattern is observed in the national accounts of wealthy countries, where the shares of health and education in expenditure have risen for decades. It is not a pathology to be cured. It is a consequence of the fact that productivity growth in some sectors makes goods cheaper than services, and the appropriate response is to ensure that the services remain accessible, which is a matter of public finance and not of technology alone.

The introduction of capable systems alters the pattern by making some tasks within the service sectors susceptible to productivity growth, and a naive reading would conclude that the cost disease will therefore be cured. A more careful reading reaches a more qualified conclusion. Services are bundles of tasks, and the tasks differ in their susceptibility to automation. Insofar as a service contains a component that is nondelegable by the principle of this book, such as the relation between teacher and student or the judgment of a physician about what to recommend to a particular patient, that component retains a fixed labor requirement, and the cost of the service cannot fall below the cost of that component. The mathematics is the same as that of Chapter 11: if a fraction of the labor in a service is irreducible, the cost reduction achievable by automating the rest is bounded by the reciprocal of that fraction. If, say, forty per cent of the labor is irreducible, the cost of the service cannot be reduced by more than a factor of two and a half, however productive the automated component becomes. The cost disease is therefore attenuated for the services affected and persists in its core.

Several consequences follow for the structure of public services. First, the nondelegable core of each service will account for a rising share of the wage bill, and of public budgets, and a society that is to afford the services it considers essential must either accept a rising share of national income devoted to them or let their quality and availability decline. The first is the better course, but it requires an institutional means of capturing the gains from productivity elsewhere in the economy and directing them to the core. Second, the professions involved in the core will change in character. The role of the physician will move toward the interpretation of automated analyses in the light of the patient's circumstances and wishes, the role of the teacher toward mentoring and the oversight of instructional systems, the role of the judge toward the explanation of decisions to those affected by them. These are roles of greater responsibility, and they require training that differs from the present training. Third, the roles that are most susceptible to displacement are those in which the work consists mainly of the tasks that can be automated, such as routine processing of claims and documents, and the persons concerned are often employed in the middle ranks of organizations, in occupations that have for decades offered stable livelihoods. The program must therefore address the transition of these persons with measures commensurate to their numbers: retraining supported by income during the retraining, placement into the expanding care and mentoring functions, and where necessary supplementary income, financed by the gains that automation produces.

The profession is not only an economic category; it is a body of knowledge, a code of conduct, and an institution of mutual accountability, and the integrity of these goods deserves protection. The licensing of professionals exists in part to ensure that persons who take responsibility for the lives and affairs of others are competent and answerable. If the profession's characteristic functions are performed by systems whose operators are not licensed, whose training is undocumented, and whose failures are not subject to professional discipline, the protection is lost. The program therefore requires that any system performing a function reserved to a profession operate under the supervision of a licensed member of that profession who is personally responsible for the result, that the record of the system's contribution be preserved, and that the standards of the profession be extended to cover the use of such systems. The professional bodies have an important role in this extension, as they do in the setting of standards of competence and ethics, and the program regards their participation as necessary and not as an obstacle.

A few words are needed about the public interest in the adoption of such systems by government. Governments already use algorithmic tools to determine eligibility for benefits, to assess risk in policing and sentencing, to detect fraud, and to allocate resources, and the record of such use includes instructive failures in which errors affected thousands of persons before they were detected, and in which the persons affected were poorly placed to contest the decisions. The principles stated in Chapter 13 bear directly on this: a decision that affects the rights or benefits of a person is a judgment in the relevant sense and must be made by a person who can be asked to explain it and who is accountable for it, and the affected person must have access to the grounds and an effective means of appeal. The use of automated tools to assist officials is compatible with these requirements; the use of them to replace the official, with the possibility of review reduced to a formality, is not. The program recommends that government adoption of such systems be preceded by public impact assessment, accompanied by publication of the system's purpose, data, and measured error rates, and followed by independent audit.

\section*{Formalization}

Two models are used, the first for the sectoral pattern of the cost disease and the second for the limit on cost reduction from partial automation.

Consider an economy with two sectors. In the progressive sector, output is $Y_p=A_p(t)L_p$ with $A_p(t)=A_0e^{gt}$, $g>0$. In the stagnant sector output is $Y_s=BL_s$ with constant $B>0$. Labor is mobile, so the wage $w$ is common, and prices equal unit labor costs: $p_p=w/A_p(t)$ and $p_s=w/B$.

\begin{proposition}
The relative price of the stagnant sector's output is $p_s/p_p=A_0e^{gt}/B$, increasing at rate $g$, and doubles every $\ln2/g$ units of time. If the economy holds the ratio of real outputs fixed at $Y_s/Y_p=\varrho>0$, then the employment share of the stagnant sector is
\[
\frac{L_s}{L}=\frac{1}{1+B\,e^{-gt}/(\varrho A_0)},
\]
which is increasing in $t$ and tends to $1$ as $t\to\infty$.
\end{proposition}

\begin{proof}
The first claim follows from the pricing rule. For the second, $Y_s/Y_p=BL_s/(A_0e^{gt}L_p)=\varrho$ implies $L_p/L_s=B/(\varrho A_0e^{gt})$, and $L_s/L=1/(1+L_p/L_s)$.
\end{proof}

With $g=0.02$ per year, the relative price rises by the factor $e^{0.02\cdot35}\approx2.01$ in thirty-five years and by $e\approx2.72$ in fifty. If instead the demand for the stagnant sector's output is such that the nominal expenditure share is fixed, the employment share rises with the relative price. In either case, the sector whose productivity does not grow absorbs a rising share of labor, and the proposition identifies what a public finance system must do to keep its services accessible.

For partial automation, let a service require, per unit of output, $\ell_J$ hours of nondelegable labor and $\ell_A$ hours of automatable labor, and let automation raise the productivity of the automatable part by the factor $1+\gamma$, $\gamma\ge0$, so that the labor requirement becomes $\ell(\gamma)=\ell_J+\ell_A/(1+\gamma)$. Define the nondelegable share $s_J=\ell_J/(\ell_J+\ell_A)$.

\begin{proposition}
The cost ratio $\ell(0)/\ell(\gamma)$ is increasing in $\gamma$ and bounded above by $1/s_J$, with $\ell(0)/\ell(\gamma)\to1/s_J$ as $\gamma\to\infty$. To attain a cost reduction by the factor $\kappa<1/s_J$ requires $1+\gamma=(1-s_J)\kappa/(1-s_J\kappa)$.
\end{proposition}

\begin{proof}
Normalize $\ell_J+\ell_A=1$. Then $\ell(\gamma)=s_J+(1-s_J)/(1+\gamma)$, which decreases in $\gamma$ to $s_J$. Setting $1/\ell(\gamma)=\kappa$ and solving for $1+\gamma$ gives the stated expression.
\end{proof}

For $s_J=0.4$ the supremum of the cost reduction is $2.5$. To attain a reduction by $\kappa=2$ requires $1+\gamma=0.6\cdot2/(1-0.8)=6$, a sixfold productivity gain in the automatable part, and to reach $\kappa=2.4$ requires $1+\gamma=0.6\cdot2.4/(1-0.96)=36$. The strongly diminishing return shows that the final portion of the reduction demands increases in productivity that are not plausibly obtained, and it is the quantitative justification for the planning assumption that public services will retain a large wage bill in their nondelegable core.
